%% kapa: exam=1 round=33 year=2022 period=1 form=A subject=civil_law subject_ko=민법 seq=34 number=34
%% source: Q-Net 감정평가사 기출문제 (cst00309) artlSeq=5209997 file=raw/1008/5209997/2022년도 제33회 감정평가사 1차시험 1교시(민법,경제학,부동산학).pdf
%% conversion: kordoc 4.24.51 -> tools/gen1.py (auto); review=proofread-sonnet; flags=none
\begin{question}{34}
민사유치권자 甲에 관한 설명으로 옳지 않은 것은?
\begin{choices}
  \item 甲이 수취한 유치물의 과실은 먼저 피담보채권의 원본에 충당하고 그 잉여가 있으면 이자에 충당한다.
  \item 甲은 피담보채권의 변제를 받기 위하여 유치물을 경매할 수 있다.
  \item 甲이 유치권을 행사하더라도 피담보채권의 소멸시효의 진행에는 영향을 미치지 않는다.
  \item 甲은 채무자의 승낙이 없더라도 유치물의 보존에 필요한 사용은 할 수 있다.
  \item 甲은 피담보채권 전부의 변제를 받을 때까지 유치물 전부에 대하여 그 권리를 행사할 수 있다.
\end{choices}
\end{question}
\answer{1}
